\documentclass[11pt]{article}
\usepackage[margin=1in]{geometry}
\usepackage{times}
\usepackage{hyperref}
\hypersetup{colorlinks=true, linkcolor=black, urlcolor=blue, citecolor=black}
\title{Rotating a Magnet Faster Does Not Rotate the Rule Book: Paramahamsa Tewari's Space Power Generator}
\author{Flyxion \\ \small Independent Researcher}
\date{}
\begin{document}
\maketitle

\begin{abstract}
Paramahamsa Tewari, a former Indian nuclear engineer, has since the 1980s promoted a "Space Power Generator" claimed to produce electrical output exceeding its mechanical input by exploiting a specific rotor and coil geometry said to draw energy directly from the surrounding vacuum or "space" medium, tying the claim to his broader theoretical framework proposing that mass and charge arise from a substratum of absolute space rather than from standard particle physics. This paper argues that the reported anomalous efficiency figures trace to conventional generator measurement errors, particularly comparing shaft mechanical input power to electrical output power without properly accounting for reduced magnetic drag torque in the specific rotor configuration used, that no version of the device has undergone a full mechanical-to-electrical energy balance test under independent instrumentation published in a peer-reviewed engineering venue, and that Tewari's broader theoretical framework, while internally elaborate, has not been shown to reduce to or extend standard electromagnetic theory in the domains where that theory is already extremely well confirmed, a minimum bar any replacement framework must clear before its novel predictions can be taken as more than speculative.
\end{abstract}

\section{Introduction}

Paramahamsa Tewari, who held a senior engineering position at India's Kaiga Atomic Power Station, has since the 1980s promoted a generator design using a specific arrangement of a rotating magnet and stationary coil geometry, claiming that certain configurations reduce the generator's back-reaction torque, the drag a conventional generator's own magnetic field exerts against the rotor as current is drawn, below the level standard electromagnetic theory predicts, allowing more electrical energy to be extracted per unit of mechanical input than conventional generators achieve, and framing this within a broader personal theoretical system describing mass and charge as vortex structures in an underlying medium called "space" or "absolute space," distinct from and, in Tewari's account, more fundamental than the electromagnetic field of standard physics.

\section{The Measurement Error This Class of Claim Commonly Involves}

Claims of generator efficiency exceeding conventional expectations most commonly trace to underestimating actual mechanical input power, particularly when a motor's electrical input, rather than its actual mechanical shaft output accounting for the motor's own conversion losses, is used as a proxy for mechanical power delivered to the generator under test, a substitution that systematically understates true mechanical input by the motor's own inefficiency, typically ten to twenty percent or more for small motors, in a way that can produce an apparent generator efficiency exceeding one hundred percent even when the generator itself performs entirely conventionally. Rigorous measurement requires direct mechanical torque and rotational speed instrumentation on the shaft connecting motor to generator, independent of either device's own electrical characteristics, a standard the publicly available documentation of Tewari's device tests has not been shown to consistently meet across independent, peer-reviewed replications.

\section{Why "Reduced Back-Reaction Torque" Is Not, by Itself, an Anomaly}

Generator designs genuinely do vary in the magnetic drag torque they produce for a given electrical output, a well studied and entirely conventional area of electrical machine design, and specific rotor and stator geometries can and do reduce certain loss mechanisms relative to a poorly optimized baseline design; this is a real engineering consideration but operates entirely within, rather than in violation of, standard electromagnetic theory's constraint that a generator's total electrical output energy cannot exceed the total mechanical input energy once all input and output are properly measured. A claim that a specific geometry reduces back-reaction torque below what standard theory predicts for that geometry, rather than merely reducing losses relative to a different, less efficient design, is the specific claim requiring extraordinary substantiation, and Tewari's published materials have not, in independent engineering review, demonstrated this distinction is what the reported efficiency figures actually reflect rather than the more mundane input-measurement substitution described above.

\section{The Broader Theoretical Framework Does Not Lower This Bar}

Tewari's "space vortex theory," proposing that particles and fields emerge from structures in an absolute space medium, is occasionally cited by supporters as independent theoretical grounding for the generator's claimed anomalous behavior. Any replacement or extension of standard electromagnetic theory must, at minimum, be shown to reproduce that theory's already extraordinarily well-confirmed predictions, from precision spectroscopy to particle accelerator physics to everyday electrical engineering, before its novel predictions in an untested regime, such as generator back-reaction torque, can be taken as more than speculative extrapolation, and no peer-reviewed demonstration that space vortex theory reproduces this existing body of confirmed electromagnetic prediction has been published in a mainstream physics venue, leaving the theoretical framework in a position where it cannot yet lend independent support to the generator claim regardless of the generator's own measurement issues.

\section{The Entropic Gravity Framing}

Tewari's space vortex theory bears a structural resemblance to other "everything emerges from an underlying medium" frameworks addressed elsewhere in this series, including the torsion field literature, in that both propose a substratum from which conventional physical quantities emerge without demonstrating that the substratum's dynamics reduce to already confirmed physics in the appropriate limits. An entropic account of gravitation, while itself proposing that spacetime curvature emerges from a more fundamental redistribution process, is held in the user's broader research program to a comparable standard, reproducing or connecting to established general-relativistic and thermodynamic results in appropriate limits, which is the specific bar this essay argues Tewari's framework has not yet cleared for electromagnetism.

\section{Objections}

Supporters point to Tewari's credentials as a career nuclear engineer as evidence the claims deserve more serious consideration than those from less credentialed inventors elsewhere in this series. Professional engineering credentials in one specialty, nuclear reactor engineering, do not substitute for independent, peer-reviewed replication of a claim in a different specialty, electrical generator design and measurement, and the specific instrumentation objection raised above applies regardless of the claimant's professional background, since a credentialed engineer's mechanical-to-electrical power measurement can suffer from the same motor-efficiency substitution error as an uncredentialed hobbyist's, if the same measurement shortcut is taken.

\section{Conclusion}

The Space Power Generator's reported efficiency gains are consistent with a well documented and common measurement substitution, using motor electrical input rather than measured mechanical shaft power as the baseline, and Tewari's broader theoretical framework has not been independently shown to reproduce standard electromagnetic theory's confirmed predictions, a precondition any candidate replacement theory must satisfy before its novel claims in an untested regime can be evaluated as more than speculative.

\end{document}
